\section{Related Work}
\label{sec:related-work}

\subsection{Agent Security}

LLM agents connect language models with tools, files, APIs, code execution,
and other external capabilities. Agent security research has extensively
studied attacks that manipulate instructions controlling agent behavior.
Prompt injection and indirect prompt injection place adversarial instructions
in user input or external content consumed by an agent
\cite{greshake2023notwhat,liu2024promptinjection,injecagent2024,
debenedetti2024agentdojo}. Tool selection attacks manipulate how agents choose
and invoke external capabilities \cite{toolhijacker2026}. Agent benchmarks
including InjecAgent, AgentDojo, and Agent Security Bench evaluate such attacks
across tools, tasks, and security policies
\cite{injecagent2024,debenedetti2024agentdojo,asb2024}.

Defense research focuses on controlling how agents interpret instructions and
invoke external capabilities. Instruction hierarchy establishes priority
relationships among trusted and untrusted instructions
\cite{wallace2024instructionhierarchy}. IPIGuard uses tool dependency graphs
to restrict actions affected by indirect prompt injection
\cite{ipiguard2025}. AgentArmor applies program analysis to agent runtime
traces \cite{wang2025agentarmor}, while AttriGuard uses counterfactual
execution to attribute tool calls to user intent and identify actions induced
by untrusted content \cite{he2026attriguard}. Agent auditing systems also study
security properties across complete agent applications and execution traces
\cite{agentaudit2026}.

Agent skills extend the attack surface to reusable third party packages that
supply instructions and executable artifacts before task execution. \system
focuses on security evidence associated with packaged skill behavior.


\subsection{Skill Security}

Agent skills combine natural language instructions, executable scripts, tool
interfaces, dependencies, and supporting artifacts. Skill security therefore
covers both instruction level manipulation and conventional software supply
chain behavior. Skill-Inject studies attacks embedded in skill files and
measures data exfiltration, destructive actions, and other harmful behavior
during agent execution \cite{schmotz2026skillinject}. SkillJect develops
automated and stealthy skill based prompt injection across instructions and
auxiliary scripts \cite{skillject2026}. Under the Hood of
\texttt{SKILL.md} studies semantic supply chain attacks that manipulate skill
discovery and selection through skill descriptions and registry behavior
\cite{underhoodskillmd2026}. Real world studies have also documented
credential theft, remote code execution, persistence, data exfiltration, and
malicious dependency delivery in public skill ecosystems
\cite{skillscan2026,liu2026maliciousskillswild}.

Several systems audit skills before deployment. SkillScan combines static
analysis with LLM based semantic classification and develops a taxonomy of
prompt injection, data exfiltration, privilege escalation, and supply chain
risks \cite{skillscan2026}. Semia represents skill interfaces and prose
semantics as structured facts and evaluates security properties through
deterministic reachability queries \cite{wen2026semia}. SkillFortify studies
formal analysis, capability verification, trust relationships, and dependency
resolution for agent skill supply chains \cite{skillfortify2026}. Industry
and open source tools, including Cisco AI Defense Skill Scanner, ClawVet, and
other skill scanners, provide rule based, semantic, and package level security
analysis \cite{ciscoSkillScanner2026,clawvet2026,
syedabbastSkillScanner2026}. SkillsMetric studies the detection boundary of
static analysis across different forms of malicious skill behavior
\cite{chen2026skillsmetric}.

Concurrent work by Ji et al. studies semantics preserving evasion against
skill scanners and proposes SkillDetonate for runtime skill analysis.
SkillDetonate executes skills in a sandbox and tracks security relevant
information flow across agent context, files, processes, and network activity
\cite{ji2026cloakdetonate}. Runtime auditing is particularly relevant to
behavior that is obfuscated, staged, or materialized only during execution.

Large scale measurement and benchmark studies provide a complementary view of
the ecosystem. Agent Skills in the Wild analyzes 31,132 unique public skills
with SkillScan and develops an empirical taxonomy of fourteen vulnerability
patterns across prompt injection, data exfiltration, privilege escalation, and
supply chain risks \cite{skillscan2026}. Malicious Agent Skills in the Wild
studies verified malicious public skills and their attack structures
\cite{liu2026maliciousskillswild}. MalSkillBench constructs a runtime verified
benchmark covering both code injection and prompt injection and evaluates
skill specific detectors, supply chain scanners, and prompt injection defenses
under a common setting \cite{guo2026malskillbench}.

Existing research covers skill vulnerability discovery, static auditing,
runtime behavior analysis, and ecosystem measurement. \system extends this
body of work with aligned instruction and runtime evidence for security
adjudication.


\subsection{Security Provenance}

System provenance reconstructs security relevant dependencies from operating
system events. Backtracking systems trace an observed intrusion to the
processes and objects that contributed to the event
\cite{king2005backtracking}. SLEUTH reconstructs attack scenarios from
commodity audit data \cite{hossain2017sleuth}, while whole system provenance
systems capture dependencies among processes, files, sockets, and other system
objects for online analysis \cite{pasquier2018camquery,
bates2015trustworthy}.

Later work applies provenance graphs to intrusion detection, investigation, and
alert reduction. NoDoze ranks provenance based alerts to reduce analyst
workload \cite{hassan2019nodoze}. UNICORN models long running system
provenance for advanced persistent threat detection
\cite{han2020unicorn}. AIQL supports attack investigation over large volumes
of system monitoring data \cite{gao2018aiql}, while ProTracer combines logging
and taint tracking to reduce provenance collection cost
\cite{ma2016protracer}. HOLMES and POIROT correlate low level audit evidence
with higher level attack behavior for investigation and threat hunting
\cite{milajerdi2019holmes,milajerdi2019poirot}. KAIROS applies provenance
analysis to practical intrusion detection and investigation over whole system
telemetry \cite{cheng2024kairos}.

Information flow analysis provides another representation of security relevant
paths. Dynamic taint analysis tracks data through program execution and has
been used for exploit analysis and privacy monitoring
\cite{newsome2005dynamic,schwartz2010all,enck2010taintdroid}. Static taint and
value flow analysis identify paths from security sensitive sources to sinks
without executing the program \cite{livshits2005static,tripp2009taj,
sui2016svf}. Code property graphs combine syntax, control flow, and data flow
relations in a common graph representation for vulnerability discovery
\cite{yama2014cpg}.

Agent skill analysis extends provenance reasoning to security relevant
behavior spanning natural language instructions and runtime execution.
\system connects evidence across both sources.